% =====================================================================
%  RESEARCH EXPERIENCE — CV section
%
%  Format follows the reference CV: bold title with right-aligned date,
%  supervisor line, a two-to-three line description, then contribution
%  bullets.
%
%  Everything not marked <<LIKE THIS>> is sourced — thesis numbers come
%  from the SafeSAC report, competition facts from IEEE SPS, ESA Phi-lab
%  and the BUET EEE co-curricular record.
%
%  TENSE: past for the thesis and VIP Cup, present for ClearSAR while it
%  is still live. Fix the tense if ClearSAR has concluded.
% =====================================================================

\section{Research Experience}

% ---------------------------------------------------------------- 1
\textbf{Undergraduate Thesis --- Safe Deep Reinforcement Learning for Vehicle-to-Grid Voltage
Support in Weak Distribution Feeders: A Physics-Aware Approach} \hfill Jul 2025 -- Jun 2026\\
Supervisor: Dr.\ Md.\ Forkan Uddin, Professor, Department of EEE, BUET

Developed \textbf{SafeSAC}, a physics-aware safe reinforcement learning controller for
bidirectional EV charging on weak radial feeders, coupling an augmented-Lagrangian Soft
Actor-Critic agent with a second-order cone safety projection rebuilt at every control step from
voltage sensitivities measured on the live network.

\begin{itemize}
  \item Formulated real-time bidirectional charging as a \textbf{constrained Markov decision process} on a modified IEEE 33-bus Baran--Wu feeder --- 95-dimensional observation, 8-dimensional continuous action, 288 five-minute steps per episode --- with explicit voltage-band and transformer-thermal constraints
  \item Designed a per-step \textbf{SOCP safety projection} (CVXPY/CLARABEL, solved in milliseconds) whose feasible set is rebuilt each step from live power-flow sensitivities, so the controller adapts to the deployment grid's R/X structure instead of assuming a nominal model
  \item Demonstrated grid-awareness quantitatively: an identical $-80$~kW V2G request is curtailed to $-54$~kW at the most sensitive station on the weak feeder while passing unaltered on a strong grid
  \item Matched the SAC-Lag baseline on voltage safety (violation rate 0.091 vs.\ 0.090, paired $t$-test $p = 0.86$) while \textbf{doubling service quality} (state-of-charge targets met 0.569 vs.\ 0.277; $+29$ percentage points, Cohen's $d = 2.04$, $p \ll 0.001$) across 25 shared-seed episodes
  \item Established robustness to train/deploy grid shift --- the unprojected baseline collapses to a non-charging policy (SoC met $= 0.000$) when moved from a strong to a weak feeder, whereas SafeSAC remains operational and delivers $+45$ percentage points more service
  \item Evaluated under a \textbf{pre-registered five-gate acceptance protocol} with paired $t$-tests, Wilcoxon signed-rank tests, bootstrap confidence intervals and Pareto analysis; released SHA-256-fingerprinted checkpoints and a complete artifact set for reproducibility
  \item Implemented the full stack in PyTorch, pandapower and CVXPY with a Gymnasium environment wrapping an in-the-loop AC power-flow solver; two-author project, <<state your half of the work in one clause>>
\end{itemize}

\vspace{6pt}

% ---------------------------------------------------------------- 2
\textbf{IEEE Signal Processing Society Video \& Image Processing (VIP) Cup 2025 --- Team EyeQ}
\hfill <<Mon>> -- Sep 2025\\
Supervisor: Dr.\ Shaikh Anowarul Fattah, Professor, Department of EEE, BUET \quad|\quad
Tutor: Anjan Kumar Bagchi

\textbf{Second Runner-Up} among four international finalist teams. Infrared--visual sensor fusion
for drone detection, tracking and payload identification in surveillance video, presented at
IEEE ICIP 2025.

\begin{itemize}
  \item Built detection and tracking models across \textbf{three modalities} --- visible (RGB), infrared, and a fused RGB--IR stream --- over roughly 45,000 pixel-registered training pairs, 6,500 validation and 13,000 test images
  \item Addressed the failure modes that motivate fusion: RGB degrades under low light, fog and glare, while thermal signatures remain detectable at night and under partial occlusion
  \item Formulated payload identification as \textbf{two-class detection} (harmful vs.\ benign) on both visible and thermal imagery, and drone--bird discrimination as a paired-modality classification problem
  \item <<YOUR SPECIFIC CONTRIBUTION --- which modality, which model, which sub-task did you personally own? A ten-person team makes this essential>>
  \item <<Any metric you can quote: mAP, precision/recall, inference latency, leaderboard position>>
\end{itemize}

\vspace{6pt}

% ---------------------------------------------------------------- 3
\textbf{ClearSAR Challenge 2026 --- European Space Agency $\Phi$-lab} \hfill <<Mon Year>> -- Present\\
<<Supervisor / Team affiliation, if any>>

Operational detection of \textbf{radio-frequency interference in Sentinel-1 synthetic aperture
radar imagery} --- an ESA $\Phi$-lab grand challenge benchmarking methods robust and scalable
enough to integrate into the mission's live processing chain. Results announced at IEEE ICIP 2026,
Tampere, Finland.

\begin{itemize}
  \item Developing automated RFI detection over a curated Sentinel-1 dataset of <<3,940>> <<quicklook / GRD>> products with annotated interference events, spanning realistic mission conditions
  \item Targeting a problem left unaddressed by standard Sentinel-1 processing, where unmitigated RFI produces false alarms and degrades every downstream Earth-observation analysis
  \item <<YOUR METHOD --- architecture, preprocessing, loss, ensembling>>
  \item <<YOUR RESULT --- leaderboard score, rank, metric>>
  \item <<YOUR ROLE if this is a team entry>>
\end{itemize}

% =====================================================================
%  RESOLVED from the final thesis (SafeSAC_Final.pdf)
%    Supervisor : Dr. Md. Forkan Uddin, Professor, Dept. of EEE, BUET
%    Co-author  : Sad Sami (ID 202006150)
%    Course     : EEE 400, July 2025 term; submitted June 2026
%    Every headline number is unchanged from the earlier draft, so the
%    bullets above are accurate against the final.
%
%  STILL NEEDED
%    1. How the thesis work split between you and Sad Sami. The thesis
%       has no per-author contribution table (unlike your nine course
%       reports), so a reviewer has no way to tell - and will ask.
%    2. Your personal contribution to VIP Cup. The EyeQ roster lists
%       ten students, so "Team EyeQ did X" is not a claim about you.
%    3. ClearSAR: which track, your role, and current status.
%
%  FIX IN THE THESIS BEFORE SENDING IT ANYWHERE
%    a. The residential load figure caption still reads "placeholder
%       parametric model (replace before publication)".
%    b. Student IDs read 202006137 / 202006150 on the thesis but
%       2006137 / 2006150 across all nine project reports.
% =====================================================================
